%% kapa: exam=1 round=34 year=2023 period=1 form=A subject=real_estate subject_ko=부동산학원론 seq=28 number=108
%% source: Q-Net 감정평가사 기출문제 (cst00309) artlSeq=5211789 file=raw/1008/5211789/2023년도 제34회 감정평가사 1차시험 1교시.pdf
%% conversion: kordoc 4.24.51 -> tools/gen1.py (auto); review=proofread-sonnet; flags=none
\begin{question}{108}
부동산신탁에 관한 설명으로 옳지 않은 것은?
\begin{choices}
  \item 신탁회사의 전문성을 통해 이해관계자들에게 안전성과 신뢰성을 제공해 줄 수 있다.
  \item 부동산신탁의 수익자란 신탁행위에 따라 신탁이익을 받는 자를 말하며, 위탁자가 지정한 제3자가 될 수도 있다.
  \item 부동산신탁계약에서의 소유권 이전은 실질적 이전이 아니라 등기부상의 형식적 소유권 이전이다.
  \item 신탁재산은 법률적으로 수탁자에게 귀속되지만 수익자를 위한 재산이므로 수탁자의 고유재산 및 위탁자의 고유재산으로부터 독립된다.
  \item 부동산담보신탁은 저당권 설정보다 소요되는 경비가 많고, 채무불이행 시 부동산 처분 절차가 복잡하다.
\end{choices}
\end{question}
\answer{5}
